\subsection{Brak podpisu}

Atak polegał na pobraniu poprawnego tokena, zamianie zawartego w nim id użytkownika na inny (przykładowo 1 dla konta admina), a następnie wysłaniu go do chatbota, który w podstawowej aplikacji odsyła odświeżone tokeny.

Przykładowy FAKE\_TOKEN był równy: 
\begin{lstlisting}
eyJ0eXAiOiAiSldUIiwgImFsZyI6ICJub25lIn0.eyJzdGF0dXMiOiAic3VjY2VzcyIsICJkYXRhIjogeyJpZCI6IDEsICJ1c2VybmFtZSI6ICJDYXJsIi
wgImVtYWlsIjogImZpbGlwQG1haWwuY29tIiwgInBhc3N3b3JkIjogImJhNmNkYWE1M2RlZDQ0Mjg1
OTYzMDE0Yzgz
ZGNjYWNmZTllMDZlOWI0Mzg3NjNiYzRiYTFhMjM4MzNlZGQ0NzciLCAicm9sZSI6ICJjdXN0b21lci
IsICJkZWx1eGVUb2tlbiI6ICIiLCAibGFzdExvZ2luSXAiOiAiMTAuNDIuMTIuMTA0IiwgInByb2Zp
bGVJbWFnZSI6ICJodHRwOi8vYXNzZXRzL3B1YmxpYy9pbWFnZXMvdXBsb2Fkcy8xLmpwZzsgc2NyaX
B0LXNyYyAndW5zYWZlLWlubGluZSciLCAidG90cFNlY3JldCI6ICIiLCAiaXNBY3RpdmUiOiB0cnVl
LCAiY3JlYXRlZEF0IjogIjIwMjUtMDUtMTIgMDk6MDQ6NTIuMTU4ICswMDowMCIsICJ1cGRhdGVkQX
QiOiAiMjAyNS0wNS0xMiAxNTo1MTowMi45MTYgKzAwOjAwIiwgImRlbGV0ZWRBdCI6IG51bGx9LCAi
aWF0IjogMTc0NzA2NTYxNH0.
\end{lstlisting}

co można rozszyfrować do danych (część nieistotnych pól została pominięta):
\begin{lstlisting}
{
  "typ": "JWT",
  "alg": "none"
}
\end{lstlisting}
\begin{lstlisting}
{
  "status": "success",
  "data": {
    "id": 1,
    "username": "Carl",
    "email": "filip@mail.com",
    "password": "ba6cdaa53ded44285963014c83dccacfe9e06e9b438763bc4ba1a23833edd477",
    "role": "customer",
    [...]
  },
  [...]
}
\end{lstlisting}

oraz brak trzeciej części, czyli podpisu (co potwierdza samotna kropka w tokenie).

\begin{table}[H]
\begin{tabular}{|l|l|l|l|l|l|}
\hline
\rowcolor[HTML]{EFEFEF} 
\multicolumn{1}{|c|}{\cellcolor[HTML]{EFEFEF}\textbf{Metoda}} & \multicolumn{1}{c|}{\cellcolor[HTML]{EFEFEF}\textbf{Końcówka}} & \multicolumn{1}{c|}{\cellcolor[HTML]{EFEFEF}\textbf{Nagłowki}} & \multicolumn{1}{c|}{\cellcolor[HTML]{EFEFEF}\textbf{Ciasteczka}} & \multicolumn{1}{c|}{\cellcolor[HTML]{EFEFEF}\textbf{Dane}}                             & \multicolumn{1}{c|}{\cellcolor[HTML]{EFEFEF}\textbf{Odpowiedź/wynik}} \\ \hline
POST                                                          & /rest/chatbot/respond                                          & \begin{tabular}[c]{@{}l@{}}Authentication: \\ Bearer FAKE\_TOKEN\end{tabular}                                  & \begin{tabular}[c]{@{}l@{}}token: \\ FAKE\_TOKEN\end{tabular}                                           & \begin{tabular}[c]{@{}l@{}}\{"action":"setname",\\"query":"Carl" \}\end{tabular} & 401 Unauthenticated user \\ \hline
\end{tabular}
\end{table}

Jak widać po odpowiedzi endpoint został zabezpieczony, ponieważ automatycznie odrzuca niepodpisany token jako Unauthenticated. Dla porównania wysłanie poprawnego tokena zwróciło odpowiedź:
\begin{lstlisting}
200 {"action":"response","body":"Nice to meet you Carl, I'm Juicy"}
\end{lstlisting}
Warto jednak zaznaczyć, że nawet dla poprawnego tokena endpoint nie zwraca już odświeżonego tokena, co jest dodatkową poprawką w stosunku do podstawowej wersji aplikacji.

Dodatkowo, wysłanie niepodpisanego tokena metodą GET /rest/user/whoami z nagłówkiem i ciasteczkiem nie zwraca danych użytkownika, więc w tym miejscu aplikacja również została poprawiona.
